% =====================================================
\chapter{Conclusion and Future Work}
\label{chap:conclusion_future}

\section{Conclusion}
This study explored the evolution of diffusion-based style transfer frameworks, beginning from foundational experiments on architectures like Vision Transformers~\cite{dosovitskiy2020vit}, U–Nets~\cite{ronneberger2015unet}, and DDPMs~\cite{ho2020ddpm}, and gradually advancing toward adaptive and perceptually guided diffusion systems. The focus was to understand how these models respond to distribution shifts caused by environmental styles such as fog, haze, snow, sandstorms, and underwater scattering.
For comparitive Analysis extreme shifts were considered and through this exploration, it was observed that early frameworks like StyleShot~\cite{gao2024styleshot} and CSGO~\cite{xing2024csgocomp}, though capable of stylization, lacked volumetric realism and semantic consistency. InstantStyle~\cite{wang2024instantstyle} emerged as a more stable and physically coherent framework, offering controllable style injection through its IP-Adapter~\cite{chen2023ipadapter} mechanism. Building upon this, the adapted InstantStyle-Plus~\cite{wang2024instantstyleplus} added spatial consistency using ControlNet~\cite{zhang2023controlnet}, while the proposed InstantStyle with CSD~\cite{somepalli2024csd} Guidance integrated perceptual feedback directly into the diffusion process. This final model demonstrated the best balance between realism, controllability, and computational efficiency, producing physically consistent, depth-aware, and semantically aligned outputs. The results highlight that realistic domain adaptation through diffusion models is best achieved not by increasing complexity, but by combining adaptive structure with perceptual understanding.

\section{Future Work}
Future work can extend this research toward more physics-aware and adaptive generative pipelines. Integrating explicit depth or scattering priors into diffusion sampling could improve the realism of environmental media like fog and haze. Developing dynamic guidance schedules that adjust style strength automatically across scenes may further enhance generalization. Testing on large-scale classification and detection tasks will confirm whether these synthetic shifts enhance real-world robustness.
